\documentclass{amsart}
\usepackage{amsmath,amssymb,amsthm,hyperref}
\usepackage{algorithm}
\usepackage{algpseudocode}

\newtheorem{theorem}{Theorem}
\newtheorem{lemma}{Lemma}
\newtheorem{proposition}{Proposition}
\theoremstyle{definition}
\newtheorem{definition}{Definition}
\newtheorem{remark}{Remark}

\DeclareMathOperator{\supp}{supp}
\DeclareMathOperator{\err}{err}
\DeclareMathOperator*{\OR}{OR}
\newcommand{\eps}{\varepsilon}
\newcommand{\X}{\{0,1\}^n}

\begin{document}

\title{Efficient semi-supervised PAC learning of two-sided disjunctions}
\maketitle

\section{Statement and answer}

We work in the model described in the problem. Examples lie in $X=\X$. A
\emph{two-sided disjunction} is a pair $h=(h^+,h^-)$ where, for some subsets
$S^+,S^-\subseteq\{1,\dots,n\}$,
\[
  h^+(x)=\OR_{i\in S^+} x_i=\mathbf 1[\,S^+\cap P(x)\neq\emptyset\,],\qquad
  h^-(x)=\OR_{i\in S^-} x_i=\mathbf 1[\,S^-\cap P(x)\neq\emptyset\,],
\]
with $P(x):=\{i:x_i=1\}$ (and the empty disjunction is identically $0$). The
\emph{classifier} associated with $h$ uses $h^+$ only: $x$ is labeled $+1$ if
$h^+(x)=1$ and $-1$ otherwise. For a distribution $D$ on $X$ and a target
$h^\ast=(h^{\ast+},h^{\ast-})$, the error of a hypothesis classifier $g$ is
\[
  \err_{D,h^\ast}(g):=\Pr_{x\sim D}\bigl[g(x)\neq h^{\ast+}(x)\bigr].
\]
A hypothesis $h$ is \emph{compatible} with $D$ if for every $x\in\supp(D)$ we
have $h^+(x)=\lnot h^-(x)$, i.e.\ exactly one of $h^+(x),h^-(x)$ equals $1$. Let
$C$ be the class of two-sided disjunctions with this compatibility notion, and
for a fixed $D$ let
\[
  C_D:=\{h\in C: h \text{ is compatible with } D\}
\]
be the compatibility-filtered subclass. Write
$m_{\mathrm{filt}}(D,\eps,\delta)$ for the
\emph{compatibility-filtered (labeled) sample complexity}: the smallest number
of labeled examples that suffices, in the standard realizable PAC sense, to
output a classifier of error at most $\eps$ with probability at least
$1-\delta$ for every target in $C_D$ labeling examples drawn from $D$.

\begin{theorem}\label{thm:main}
The answer to the problem is \textbf{yes}. There is an algorithm
$\mathcal A$ for the semi-supervised PAC model of $C$ such that:
\begin{enumerate}
\item[(a)] $\mathcal A$ runs in time polynomial in $n$, $1/\eps$ and
$1/\delta$;
\item[(b)] on every unknown distribution $D$ and every target
$h^\ast\in C_D$, with probability at least $1-\delta$ the output classifier
has error at most $\eps$, and the number of labeled examples used is
\[
  m \;=\; \Bigl\lceil \tfrac{1}{\eps}\bigl(n\ln 2+\ln(1/\delta)\bigr)\Bigr\rceil
  \;\le\; p\bigl(n,1/\eps,\log(1/\delta)\bigr)\cdot
          \bigl(m_{\mathrm{filt}}(D,\eps,\delta)+1\bigr)
\]
for a fixed polynomial $p$;
\end{enumerate}
while using a number of unlabeled examples that is polynomial in $n$, $1/\eps$
and $\log(1/\delta)$ (indeed, the algorithm needs no unlabeled examples, so in
particular it uses a polynomial number of them).
\end{theorem}

The remainder of the paper proves Theorem~\ref{thm:main}. The structure is:
Section~\ref{sec:alg} gives the explicit algorithm and proves it is a correct,
polynomial-time PAC learner using $m=\lceil\frac1\eps(n\ln2+\ln(1/\delta))\rceil$
labeled examples (Lemma~\ref{lem:upper}); Section~\ref{sec:lower} proves the
quantitative comparison to $m_{\mathrm{filt}}$ (Lemma~\ref{lem:compare}); the
two combine to give the theorem.

\section{The algorithm and its labeled-sample upper bound}\label{sec:alg}

The classifier produced by a two-sided disjunction depends only on $h^+$, which
is a monotone disjunction. We therefore learn $h^{\ast+}$ directly with the
classical elimination (list-and-cross-off) algorithm. The compatibility
constraint and the unlabeled sample are \emph{not} needed for correctness; they
are only relevant to the comparison in Section~\ref{sec:lower}.

\begin{algorithm}[h]
\caption{Elimination learner for the classifier of $C$}\label{alg:elim}
\begin{algorithmic}[1]
  \Require parameters $\eps,\delta\in(0,1)$; oracle access to labeled examples
           $(x,y)$, $x\sim D$, $y=h^{\ast+}(x)$.
  \Ensure a classifier $g$ which is a monotone disjunction over $\{1,\dots,n\}$.
  \State $m \gets \bigl\lceil \tfrac1\eps\bigl(n\ln 2+\ln(1/\delta)\bigr)\bigr\rceil$
  \State $T \gets \{1,2,\dots,n\}$
         \Comment{$T$ = set of variables still allowed in $g$}
  \For{$t=1$ \textbf{to} $m$}
     \State draw a labeled example $(x,y)$
     \If{$y=-1$} \Comment{$h^{\ast+}(x)=0$}
        \State $T \gets T \setminus P(x)$
            \Comment{remove every variable set to $1$ in $x$}
     \EndIf
  \EndFor
  \State \Return the classifier $g(x):=\OR_{i\in T} x_i$
         \big(equivalently $g(x)=\mathbf 1[\,T\cap P(x)\neq\emptyset\,]$\big)
\end{algorithmic}
\end{algorithm}

\begin{lemma}\label{lem:upper}
Fix any $D$ on $X$ and any target $h^\ast\in C$ (compatibility is not used).
On a labeled sample of size
$m=\lceil\frac1\eps(n\ln2+\ln(1/\delta))\rceil$, Algorithm~\ref{alg:elim}:
\begin{enumerate}
\item runs in time $O(mn)=\mathrm{poly}(n,1/\eps,\log(1/\delta))$ and uses $m$
labeled examples and no unlabeled examples;
\item outputs a classifier $g$ with
$\Pr[\err_{D,h^\ast}(g)\le\eps]\ge 1-\delta$,
the probability over the labeled sample.
\end{enumerate}
\end{lemma}

\begin{proof}
\emph{Running time and sample size.} The loop runs $m$ times; each iteration
draws one labeled example, inspects the at most $n$ coordinates of $x$ and
deletes the corresponding elements from $T$, costing $O(n)$. Output is $O(n)$.
Total time $O(mn)$. With
$m=\lceil\frac1\eps(n\ln2+\ln(1/\delta))\rceil$ this is polynomial in
$n,1/\eps,\log(1/\delta)$, hence in $n,1/\eps,1/\delta$. No unlabeled example is
ever requested.

\emph{Structure of the output.} Let $S^\ast:=S^{\ast+}$ be the index set of the
target classifier $h^{\ast+}=\OR_{i\in S^\ast}x_i$, and let $T_{\mathrm{out}}$ be
the final value of $T$. We first record two invariants.

\textbf{(I1) $S^\ast\subseteq T_{\mathrm{out}}$.}
A variable $i$ is deleted from $T$ only when, on some example $(x,-1)$, we have
$i\in P(x)$, i.e.\ $x_i=1$. If $i\in S^\ast$ then $x_i=1$ forces
$h^{\ast+}(x)=\OR_{j\in S^\ast}x_j=1$, so the label of $x$ would be $+1$, not
$-1$. Hence no $i\in S^\ast$ is ever deleted, and $S^\ast\subseteq
T_{\mathrm{out}}$.

\textbf{(I2) $g$ is consistent with the whole sample.} Consider any sampled
example $(x,y)$.
\begin{itemize}
\item If $y=+1$, then $S^\ast\cap P(x)\neq\emptyset$; by (I1)
$S^\ast\subseteq T_{\mathrm{out}}$, so $T_{\mathrm{out}}\cap P(x)\neq\emptyset$
and $g(x)=1=y$.
\item If $y=-1$, the iteration that processed $(x,-1)$ deleted all of $P(x)$
from $T$, and elements are never re-added, so
$T_{\mathrm{out}}\cap P(x)=\emptyset$ and $g(x)=0$, i.e.\ $g$ labels $x$ as
$-1=y$.
\end{itemize}
Thus $g$ agrees with the target on every sampled example.

\textbf{(I3) one-sided dominance.} For every $x\in X$, since
$S^\ast\subseteq T_{\mathrm{out}}$ we have
$S^\ast\cap P(x)\subseteq T_{\mathrm{out}}\cap P(x)$, whence
$h^{\ast+}(x)=1\Rightarrow g(x)=1$. So $g$ never produces a false negative:
$\{x:g(x)\neq h^{\ast+}(x)\}=\{x:g(x)=1,\ h^{\ast+}(x)=0\}$.

\emph{Generalization (cardinality/Occam bound, realizable case).}
Let $\mathcal H$ be the class of all monotone disjunctions
$g_T(x)=\OR_{i\in T}x_i$, $T\subseteq\{1,\dots,n\}$; then
$|\mathcal H|=2^n$ and the output $g=g_{T_{\mathrm{out}}}\in\mathcal H$. The
target classifier $h^{\ast+}=g_{S^\ast}\in\mathcal H$, so the learning problem
is realizable. We use the standard finite-class realizable PAC bound, proved
here for completeness. Call $g_T\in\mathcal H$ \emph{bad} if
$\err_{D,h^\ast}(g_T)>\eps$. For a fixed bad $g_T$, the probability that it
agrees with the target on all $m$ i.i.d.\ sampled examples is
\[
  \Pr\bigl[g_T\text{ consistent with the sample}\bigr]
   =\bigl(1-\err_{D,h^\ast}(g_T)\bigr)^m
   <(1-\eps)^m\le e^{-\eps m}.
\]
By a union bound over the at most $|\mathcal H|=2^n$ bad hypotheses,
\[
  \Pr\bigl[\exists\text{ bad } g_T\text{ consistent with the sample}\bigr]
   \le 2^n e^{-\eps m}.
\]
With $m=\lceil\frac1\eps(n\ln2+\ln(1/\delta))\rceil\ge
\frac1\eps(n\ln2+\ln(1/\delta))$ we get
$\eps m\ge n\ln2+\ln(1/\delta)$, hence
$2^n e^{-\eps m}=e^{n\ln2-\eps m}\le e^{-\ln(1/\delta)}=\delta$. By (I2) the
output $g$ \emph{is} consistent with the sample; therefore, with probability at
least $1-\delta$, $g$ is not bad, i.e.\ $\err_{D,h^\ast}(g)\le\eps$.
\end{proof}

\section{Comparison with the filtered sample complexity}\label{sec:lower}

We now show that the labeled count $m$ of Lemma~\ref{lem:upper} matches the
compatibility-filtered sample complexity up to a polynomial factor, as required
in clause (b). The point is purely quantitative: $m$ is already a polynomial in
$n,1/\eps,\log(1/\delta)$, and the filtered complexity is at least a constant
whenever it is not vacuous; the degenerate (vacuous) case is absorbed by the
additive ``$+1$''.

\begin{lemma}\label{lem:compare}
Let $p(n,1/\eps,\log(1/\delta)):=\frac1\eps\bigl(n+\log_2(1/\delta)+1\bigr)$.
For every distribution $D$ and all $\eps,\delta\in(0,1)$,
\[
  m=\Bigl\lceil \tfrac1\eps\bigl(n\ln2+\ln(1/\delta)\bigr)\Bigr\rceil
  \;\le\; p\bigl(n,1/\eps,\log(1/\delta)\bigr)\cdot
          \bigl(m_{\mathrm{filt}}(D,\eps,\delta)+1\bigr).
\]
\end{lemma}

\begin{proof}
First bound $m$ from above. Using $\lceil a\rceil\le a+1$, $\ln 2<1$ and
$\ln(1/\delta)=\ln 2\cdot\log_2(1/\delta)<\log_2(1/\delta)$,
\[
  m \le \tfrac1\eps\bigl(n\ln 2+\ln(1/\delta)\bigr)+1
     \le \tfrac1\eps\bigl(n+\log_2(1/\delta)\bigr)+1
     \le \tfrac1\eps\bigl(n+\log_2(1/\delta)+1\bigr)
     = p\bigl(n,1/\eps,\log(1/\delta)\bigr),
\]
where in the third inequality we used $1\le \frac1\eps$ (as $\eps<1$). Hence
\begin{equation}\label{eq:mpoly}
  m \le p\bigl(n,1/\eps,\log(1/\delta)\bigr).
\end{equation}
Since $m_{\mathrm{filt}}(D,\eps,\delta)\ge 0$ always, we have
$m_{\mathrm{filt}}(D,\eps,\delta)+1\ge 1$, and multiplying the right-hand side of
\eqref{eq:mpoly} by this quantity only increases it:
\[
  m \;\le\; p\bigl(n,1/\eps,\log(1/\delta)\bigr)
        \;\le\; p\bigl(n,1/\eps,\log(1/\delta)\bigr)\cdot
                \bigl(m_{\mathrm{filt}}(D,\eps,\delta)+1\bigr).\qedhere
\]
\end{proof}

\begin{remark}[On the role of the ``$+1$'' and of $m_{\mathrm{filt}}$]
The additive ``$+1$'' only matters in the degenerate regime
$m_{\mathrm{filt}}(D,\eps,\delta)=0$, i.e.\ when a single fixed classifier is
already $\eps$-good for \emph{every} compatible target on $D$ (for instance when
$\supp(D)$ is a single point, so all classifiers agreeing there cannot be
distinguished and zero labeled examples already $\eps$-cover $C_D$). In every
non-degenerate case $m_{\mathrm{filt}}(D,\eps,\delta)\ge 1$ and the bound
becomes the clean statement
$m\le p(n,1/\eps,\log(1/\delta))\cdot m_{\mathrm{filt}}(D,\eps,\delta)$.
Indeed, a standard PAC lower bound (Ehrenfeucht--Haussler--Kearns--Valiant)
shows that whenever $C_D$ contains two classifiers at $D$-distance in
$(\eps,1]$ one in fact needs
$m_{\mathrm{filt}}(D,\eps,\delta)=\Omega\bigl(\frac1\eps\log\frac1\delta\bigr)$
labeled examples, so the polynomial factor $p$ is essentially the only loss.
\end{remark}

\section{Proof of Theorem~\ref{thm:main}}

Take $\mathcal A$ to be Algorithm~\ref{alg:elim} with parameters $\eps,\delta$.

Clause (a): by Lemma~\ref{lem:upper}(1) the running time is $O(mn)$ with
$m=\lceil\frac1\eps(n\ln2+\ln(1/\delta))\rceil$, which is polynomial in
$n,1/\eps$ and (since $\log(1/\delta)\le 1/\delta$) in $1/\delta$; thus $\mathcal
A$ runs in time polynomial in $n,1/\eps,1/\delta$.

Correctness and clause (b): let $D$ be any distribution and
$h^\ast\in C_D\subseteq C$ any target. By Lemma~\ref{lem:upper}(2),
$\mathcal A$ outputs a classifier $g$ with
$\Pr[\err_{D,h^\ast}(g)\le\eps]\ge1-\delta$, using exactly $m$ labeled examples
and no unlabeled examples. (Lemma~\ref{lem:upper} requires only $h^\ast\in C$,
which holds since $C_D\subseteq C$; the compatibility of $h^\ast$ with $D$ is a
strictly stronger hypothesis and is not needed.) By Lemma~\ref{lem:compare}
this labeled count satisfies
\[
  m \le p\bigl(n,1/\eps,\log(1/\delta)\bigr)\cdot
        \bigl(m_{\mathrm{filt}}(D,\eps,\delta)+1\bigr),
\]
with $p(n,1/\eps,\log(1/\delta))=\frac1\eps(n+\log_2(1/\delta)+1)$ a fixed
polynomial in $n,1/\eps,\log(1/\delta)$. Thus the labeled-sample size matches
the compatibility-filtered sample complexity up to a factor polynomial in
$n,1/\eps,\log(1/\delta)$.

Unlabeled sample: $\mathcal A$ uses $0$ unlabeled examples, which is trivially a
polynomial number.

All three requirements hold, completing the proof. \qed

\begin{remark}[Why the answer is positive ``for free'']
The semantic content of the result is that, for this concept class, unlabeled
data is not \emph{necessary} to achieve the polynomially-matched labeled
complexity: because the classifier of a two-sided disjunction is a single
monotone disjunction, the \emph{unfiltered} labeled sample complexity of $C$ is
already $O\bigl(\frac1\eps(n+\log\frac1\delta)\bigr)=
\mathrm{poly}(n,1/\eps,\log(1/\delta))$, and the filtered complexity, being the
complexity of the subclass $C_D\subseteq C$, can only be smaller. Matching the
smaller quantity up to a polynomial factor is therefore automatic once one has
an efficient learner attaining the polynomial unfiltered bound, which the
elimination algorithm provides. The compatibility structure and the unlabeled
sample, which in general (e.g.\ in the Balcan--Blum framework) are what make
matching possible, are here not required.
\end{remark}

\end{document}
